\specialchapter{引言}

本书所处理的问题，是在代数数论以及（稍后的）代数几何的发展过程中产生的（参看历史注记）。从十九世纪起，这两个理论开始显示出显著的类似性；解决它们所提出的问题的尝试，导致分离出若干一般概念，其应用范围并不限于代数数环或代数函数环；而且，一如既往，以最一般的形式来考虑它们是有利的，这样才能看清它们的真正意义以及一项研究对另一项研究的影响。本书所处理的概念原则上可应用于所有交换环及其上的模；然而必须指出，实质性的结果常常只是在某些有限性假设（这些假设在经典情形总成立）之下才得到，例如假定模是有限生成的或环是 Noether 环。

开头几章的中心概念如下：

I. 局部化与整体化。例如，让我们从一个丢番图方程组出发：

\[
\mathrm{P} _ {i} (x _ {1}, \dots , x _ {m}) = 0 \quad (1 \leqslant i \leqslant n)\tag{*}
\]

其中 \(P_{i}\) 是整系数多项式，而要找的由有理整数组成的解 \((x_{i})\) 。可以先从寻找由有理数组成的解入手来接近这个问题，这相当于把 \(P_{i}\) 的系数视为 Z 的分式域 Q 的元素、并把所要的解取在 Q 中来考虑同一问题。第二步是考察：给定一个素数 p，是否存在分母不被 p 整除的有理数解（整数解显然满足这一条件）；在这种情形下，这相当于位于 Q 中由这种形式的有理数组成的子环 \(\mathbf{Z}_{(p)}\) 之内，该子环称为 Z 对应于素数 p 的局部环。显然，从 Z 到 Q 的过渡与从 Z 到 \(\mathbf{Z}_{(p)}\) 的过渡具有同一形式：在两种情形下，所允许的分母都不属于某个素理想（分别是理想 (0) 与理想 (p)）。“局部环”这一名称也出现在代数几何中，在那里这个概念以更自然的方式出现：例如，对于复系数单变量多项式的环 C{[}X{]}，对应于素理想 (X - \(\alpha\) ) 的局部环是在点 \(\alpha\) 处“正则”的（即在该点没有极点的）有理分式环。

每一个丢番图问题，更一般地，每一个关于 A-模（A 为交换环）的问题，都可以分解为两个辅助问题：先在与 A 的各个素理想 p 相应的局部环 \(A_{p}\) 中求解（“局部化”），然后问：能否由“局部化”问题对所有 p 都有解推出最初所提的问题有解（“从局部到整体的过渡”）。第二章专门研究这一双重过程，在那里还会看到，“局部化”并不只与素理想相关，而是有更广的适用范围。

\begin{enumerate}
\def\labelenumi{\Roman{enumi}.}
\setcounter{enumi}{1}
\tightlist
\item
  局部环的完备化。局部环 A 与域同样具有只有一个极大理想 m 的性质。利用这一事实，可以通过转向商环 A/m（因后者是域）在一定程度上把关于 A-模的问题变换为关于向量空间的类似问题。如果回到丢番图方程组 (*)，这一思想无非就是“模 p 约化”原理，即把方程变成模 p 的同余式，它从数论最早的工作起就自然地出现了。
\end{enumerate}

尽管如此，我们显然不能指望以这种方式得到原问题的完整结果，而且人们很快便认识到，为得到更精确的信息，必须不仅考虑模 m 的同余，还要考虑模 \(m^{n}\) 的“更高阶”同余，其中 n 为任意整数 n \textgreater{} 0。这样就会发现，n 越大，原问题在某种意义上就被“逼近”得越紧（例如在 A = Z 的情形，理由是一个 \(\neq 0\) 的整数不可能被一个给定素数 p 的所有幂 \(p^{n}\) 整除；因此只要 n 取得足够大，这个数就会在模 \(p^{n}\) 约化中显示出它的存在）。这一思想的数学转译，是在 A 上考虑一个环拓扑（参看《一般拓扑学》第三章 §6），其中诸 \(m^{n}\) 构成 0 的一个基本邻域系。但是，当我们例如已经解决了同余方程组

\[
\mathrm{P} _ {i} (x _ {1}, \dots , x _ {m}) \equiv 0 (\mathrm{mod.} p ^ {k}) \quad (1 \leqslant i \leqslant n)\tag{**}
\]

对每个整数 \(k > 0\) 成立时，仍然不能由此推出方程组 (*) 在局部环 \(\mathbf{Z}_{(p)}\) 中有解；上述假设可以解释为：(*) 在 \(\hat{\mathbf{Z}}_{(p)}\) （即拓扑环 \(\mathbf{Z}_{(p)}\) 的完备化）中有一个解。

这样削弱后的原问题，最终被变换为关于 A/m \(^{n}\) 型局部环的类似问题；这些环由于具有幂零根，
也比一般环更接近于域；在经典代数几何中，这对应于在给定点的邻域中对问题作“微分”式的研究。

第三章以一般的方式处理拓扑概念对局部环理论的这些应用。第六章研究其中一个特殊方面，它一方面适合于代数几何的更细致的研究，而尤其适合于代数数域的算术，在那里遇到的局部环（例如 \(\mathbf{Z}_{(p)}\) ）属于一类特别简单的环，即“赋值环”类；在赋值环中，可除性是主理想集合的一个全序（参看《代数学》第六章 §1）。

对从环 A 过渡到局部环 \(A_{p}\) 或完备化 \(\hat{A}\) 的研究，揭示了这两种运算的一个共同特征，即 A-模 \(A_{p}\) 与 \(\hat{A}\) 的平坦性；除其他作用外，它使得这类 A-模与任意 A-模的张量积的使用在某种程度上类似于向量空间张量积的使用，也就是说，无需一般情形下围绕其使用的种种预防措施。与这一概念相联系的性质也适用于非交换环上的模，它们是第一章的研究对象。

\begin{enumerate}
\def\labelenumi{\Roman{enumi}.}
\setcounter{enumi}{2}
\tightlist
\item
  整数与理想的分解。对代数数域中可除性的研究，从一开始就需要在这种域 K 中引入整数的概念，以推广域 Q 中有理整数的概念。“代数整数”这一概念的一般理论（如将看到的，它与非常严格的有限性条件相联系）在第五章中展开；它可以应用于所有交换环，并且不仅在算术中意义重大，在代数几何中、甚至在场 C 上的“解析空间”的现代理论中也意义重大。
\end{enumerate}

把经典算术扩张到代数整数环的主要障碍之一，长期以来在于有理整数分解为素因子的经典分解一般不能推广到这些环。为克服这一困难，必须创建理想的理论：这样，所寻求的唯一分解对理想成立，素理想的概念取代了素数的概念。此外，这一结果可以看作“从局部到整体的过渡”得以圆满实现的典型情形：对 \(x \in K\) ，知道 K 上所有“赋值”在 x 处的值，就能在相差一个可逆整数乘法的意义下确定 x。

在比代数整数环不那么简单的环中（甚至在多个未定元的多项式环中），这一结果不再成立。然而，可以用典范的方式使每个理想对应于一个完全确定的素理想集合：在代数几何中，如果考虑 \(K^{n}\)（K 为任意交换域）中由多项式方程组 \(P_{\alpha}=0\) 定义的子簇，则该子簇的不可约分支与由诸 \(P_{\alpha}\) 生成的理想所对应的素理想集合的极小元素一一对应。此外，（若限于 Noether 环）可以对每个理想给出一个比素理想之积的分解不那么精确的“分解”：这里乘积实际上被交集取代，而素理想的幂被与所论理想的伴随素理想相联系的“准素”理想所取代（但它们并不是素理想幂的直接推广）。与理想相伴随的素理想的引入及其性质的研究是第四章的主题；那里还证明了刚才提到的“准素分解”的存在性以及某些唯一性性质；但目前看来，这些分解在应用中通常只起辅助作用，本质性的概念是与理想相伴随的素理想。

在第七章中，我们将更详细地考察这样一些环：就分解为素理想之积而言，它们的性质最接近代数整数环的性质；除其他事项外，在这些环中可以引入“除子”的概念，它是这一分解的几何侧面，在代数几何中起重要作用。

最后，第八章及以后各章将处理在代数几何中比在算术中（在那里它们成为平凡的）更受关注的概念，特别是维数的概念。

有了这些概念，我们就到达了严格意义上的代数几何的边界，这是一条不断移动且难以划出的边界。因为，如果交换代数是使代数几何在其全部一般性中得到发展的本质工具，那么反过来（如上面已经看到的），几何语言对于表述交换代数的定理、并提示一种抽象代数中自然缺乏的直觉，显得非常方便；随着把代数几何的界限不断扩大得越来越多的趋势，代数的语言与几何的语言比以往任何时候都更加趋于融合。

第一章(*)

